% Общие элементы бланков ВКР (по образцу VKR_Morozov.pdf)

% Текст шапки министерства, 6 строк (как у Морозова на задании и плане)
\newcommand{\vkrministrytext}{%
Министерство науки и высшего образования Российской Федерации\\
Федеральное государственное автономное образовательное учреждение\\
высшего образования\\
<<Московский государственный технический университет имени Н.Э. Баумана\\
(национальный исследовательский университет)>>\\
(МГТУ им. Н.Э. Баумана)}

% Двойная линия под шапкой (толстая + тонкая), во всю ширину текста
\newcommand{\vkrheaderrule}{%
  \par\vspace{1pt}\nointerlineskip
  \noindent\rule{\textwidth}{2.8pt}\par\nointerlineskip
  \vspace{1.6pt}\nointerlineskip
  \noindent\rule{\textwidth}{0.5pt}\par}

% Центрированная жирная шапка министерства + двойная линия (задание, план)
\newcommand{\vkrformministry}{%
  {\centering\fontsize{11.5pt}{12pt}\selectfont\bfseries\vkrministrytext\par}%
  \vkrheaderrule}

% Блок УТВЕРЖДАЮ — единый для задания и кал. плана (по образцу Морозова).
% Все три содержательные строки начинаются под буквой «З» (одинаковый левый
% край) и заканчиваются на одной вертикали; линии для заполнения тянутся до
% правого края блока (\hrulefill ~ подчёркивание). Подписи (Индекс)/(И.О.Фамилия)
% мелким шрифтом, прижаты вправо, минимальный интервал к линиям.
\newlength{\vkrutvw}\setlength{\vkrutvw}{5.7cm}
\newcommand{\vkrutvline}{\leavevmode\hrulefill}
% Мелкая подпись, прижатая к линии сверху (как у Морозова — почти вплотную)
\newcommand{\vkrutvcap}[1]{\par\nobreak\vskip-0.6ex\makebox[\vkrutvw][r]{\scriptsize#1}}
\newcommand{\vkrapproveinner}{%
\begin{minipage}[t]{\vkrutvw}%
\setlength{\parindent}{0pt}%
\setlength{\parskip}{0pt}%
\centering УТВЕРЖДАЮ\par
\vskip0.5em
{\raggedright\baselineskip=2.2ex
Заведующий кафедрой\,\vkrutvline
\vkrutvcap{(Индекс)}\par
\vskip0.35em
\vkrutvline
\vkrutvcap{(И.О.Фамилия)}\par
\vskip0.3em
<<\,\underline{\makebox[1.0cm]{}}\,>>\,\underline{\makebox[2.3cm]{}}\,20\,\underline{\makebox[0.8cm]{}}\,г.\par}%
\end{minipage}}

\newcommand{\vkrformapprove}{%
\noindent\hfill\vkrapproveinner}

\newcommand{\vkrtopic}{Компилятор базисного подмножества Рефала-5 с представлением данных, использующим раскрученные очереди}

% ФАКУЛЬТЕТ / КАФЕДРА / ГРУППА (значения на одной вертикали)
\newcommand{\vkrformfaculty}[3]{%
{\bfseries\fontsize{13pt}{15pt}\selectfont
\makebox[3.25cm][l]{ФАКУЛЬТЕТ}\underline{\makebox[2.35cm][c]{#1}}\\[0.18em]
\makebox[3.25cm][l]{КАФЕДРА}\underline{\makebox[2.35cm][c]{#2}}\\[0.18em]
\makebox[3.25cm][l]{ГРУППА}\underline{\makebox[2.35cm][c]{#3}}%
}}

% Строка ФАКУЛЬТЕТ/КАФЕДРА на титуле: метка + название по центру на подчёркивании
\newcommand{\vkrfacultyline}[2]{%
  \noindent\makebox[2.9cm][l]{#1}\underline{\makebox[\dimexpr\textwidth-2.9cm\relax][c]{#2}}\par}
